\chapter{Referencias}
\markboth{REFERENCIAS}{}

\begin{thebibliography}{99}
\small

\bibitem{abrs2021} Auclert, A., Bardóczy, B., Rognlie, M. y Straub, L. (2021).
``Using the Sequence-Space Jacobian to Solve and Estimate Heterogeneous-Agent
Models''. \emph{Econometrica}, 89(5), 2375--2408.
Repositorio: \url{https://github.com/shade-econ/sequence-jacobian}.

\bibitem{ars2018} Auclert, A., Rognlie, M. y Straub, L. (2018).
``The Intertemporal Keynesian Cross''. NBER Working Paper 25020.

\bibitem{ahn2018} Ahn, S., Kaplan, G., Moll, B., Winberry, T. y Wolf, C. (2018).
``When Inequality Matters for Macro and Macro Matters for Inequality''.
\emph{NBER Macroeconomics Annual}, 32, 1--75.

\bibitem{afonso2015} Afonso, G. y Lagos, R. (2015).
``Trade Dynamics in the Market for Federal Funds''. \emph{Econometrica}, 83(1), 263--313.

\bibitem{bayer2020} Bayer, C. y Luetticke, R. (2020).
``Solving Discrete Time Heterogeneous Agent Models with Aggregate Risk and Many
Idiosyncratic States by Perturbation''. \emph{Quantitative Economics}, 11(4), 1253--1288.

\bibitem{bianchibigio2022} Bianchi, J. y Bigio, S. (2022).
``Banks, Liquidity Management, and Monetary Policy''. \emph{Econometrica}, 90(1), 391--454.

\bibitem{bkm2018} Boppart, T., Krusell, P. y Mitman, K. (2018).
``Exploiting MIT Shocks in Heterogeneous-Agent Economies: The Impulse Response as a
Numerical Derivative''. \emph{Journal of Economic Dynamics and Control}, 89, 68--92.

\bibitem{carroll2006} Carroll, C. D. (2006).
``The Method of Endogenous Gridpoints for Solving Dynamic Stochastic Optimization
Problems''. \emph{Economics Letters}, 91(3), 312--320.

\bibitem{dd1983} Diamond, D. W. y Dybvig, P. H. (1983).
``Bank Runs, Deposit Insurance, and Liquidity''.
\emph{Journal of Political Economy}, 91(3), 401--419.

\bibitem{gertlerkaradi2011} Gertler, M. y Karadi, P. (2011).
``A Model of Unconventional Monetary Policy''.
\emph{Journal of Monetary Economics}, 58(1), 17--34.

\bibitem{gertlerkiyotaki2010} Gertler, M. y Kiyotaki, N. (2010).
``Financial Intermediation and Credit Policy in Business Cycle Analysis''.
En B. Friedman y M. Woodford (eds.), \emph{Handbook of Monetary Economics}, vol. 3A.

\bibitem{hs2015} Herbst, E. y Schorfheide, F. (2015).
\emph{Bayesian Estimation of DSGE Models}. Princeton University Press.

\bibitem{kmv2018} Kaplan, G., Moll, B. y Violante, G. L. (2018).
``Monetary Policy According to HANK''. \emph{American Economic Review}, 108(3), 697--743.

\bibitem{ks1998} Krusell, P. y Smith, A. A. (1998).
``Income and Wealth Heterogeneity in the Macroeconomy''.
\emph{Journal of Political Economy}, 106(5), 867--896.

\bibitem{reiter2009} Reiter, M. (2009).
``Solving Heterogeneous-Agent Models by Projection and Perturbation''.
\emph{Journal of Economic Dynamics and Control}, 33(3), 649--665.

\bibitem{rouwenhorst1995} Rouwenhorst, K. G. (1995).
``Asset Pricing Implications of Equilibrium Business Cycle Models''.
En T. F. Cooley (ed.), \emph{Frontiers of Business Cycle Research}. Princeton University Press.

\bibitem{sw2007} Smets, F. y Wouters, R. (2007).
``Shocks and Frictions in US Business Cycles: A Bayesian DSGE Approach''.
\emph{American Economic Review}, 97(3), 586--606.

\bibitem{winberry2018} Winberry, T. (2018).
``A Method for Solving and Estimating Heterogeneous Agent Macro Models''.
\emph{Quantitative Economics}, 9(3), 1123--1151.

\bibitem{young2010} Young, E. R. (2010).
``Solving the Incomplete Markets Model with Aggregate Uncertainty Using the
Krusell--Smith Algorithm and Non-Stochastic Simulations''.
\emph{Journal of Economic Dynamics and Control}, 34(1), 36--41.

\end{thebibliography}
